\documentclass[11pt,a4paper]{article}
\usepackage[utf8]{inputenc}
\usepackage[T1]{fontenc}
\usepackage{lmodern}
\usepackage[french]{babel}
\usepackage{amsmath,amssymb,mathtools}
\usepackage{icomma}
\usepackage{xcolor}
\usepackage{colortbl}
\usepackage{tikz}
\usepackage[most]{tcolorbox}
\usepackage{enumitem}
\usepackage{array}
\usepackage{booktabs}
\usepackage{fancyhdr}
\usepackage[hidelinks]{hyperref}
\usetikzlibrary{arrows.meta,positioning,calc,patterns,backgrounds,shapes.geometric,decorations.pathreplacing,decorations.markings}
\usepackage[a4paper,margin=2.15cm,headheight=26pt,headsep=12pt]{geometry}
\usepackage{titlesec}

% ---------------- Palette ----------------
\definecolor{primary}{RGB}{0,151,169}
\definecolor{primarydark}{RGB}{0,88,101}
\definecolor{defcol}{RGB}{0,114,178}
\definecolor{thmcol}{RGB}{0,140,120}
\definecolor{formulacol}{RGB}{198,93,9}
\definecolor{excol}{RGB}{0,135,90}
\definecolor{attncol}{RGB}{200,40,55}
\definecolor{methcol}{RGB}{90,90,100}
\definecolor{remarkcol}{RGB}{120,94,200}
\definecolor{barcol}{RGB}{0,151,169}
\definecolor{barcold}{RGB}{0,110,124}
\definecolor{modecol}{RGB}{200,55,55}
\definecolor{meancol}{RGB}{0,90,200}
\definecolor{medcol}{RGB}{160,90,190}
\definecolor{quartcol}{RGB}{0,135,90}
\definecolor{ogivecol}{RGB}{176,74,0}

% ---------------- Boîtes ----------------
\tcbset{boxbase/.style={enhanced,breakable,boxrule=0.9pt,arc=2.5pt,
  left=3mm,right=3mm,top=2.3mm,bottom=2.3mm,
  fonttitle=\bfseries,coltitle=white,toptitle=0.6mm,bottomtitle=0.6mm}}

\newtcolorbox{definition}[1][]{boxbase,colback=defcol!6,colframe=defcol,
  title={\textsc{Définition}\ifx\relax#1\relax\relax\else\textmd{ : \itshape #1}\fi}}
\newtcolorbox{propriete}[1][]{boxbase,colback=thmcol!7,colframe=thmcol,
  title={\textsc{Propriété}\ifx\relax#1\relax\relax\else\textmd{ : \itshape #1}\fi}}
\newtcolorbox{formule}[1][]{boxbase,colback=formulacol!7,colframe=formulacol,
  title={\textsc{Formule}\ifx\relax#1\relax\relax\else\textmd{ : \itshape #1}\fi}}
\newtcolorbox{exemple}[1][]{boxbase,colback=excol!6,colframe=excol,
  title={\textsc{Exemple}\ifx\relax#1\relax\relax\else\textmd{ : \itshape #1}\fi}}
\newtcolorbox{methode}[1][]{boxbase,colback=methcol!8,colframe=methcol,sharp corners,
  title={\textsc{Méthode}\ifx\relax#1\relax\relax\else\textmd{ : \itshape #1}\fi}}
\newtcolorbox{remarque}[1][]{boxbase,colback=remarkcol!7,colframe=remarkcol,
  title={\textsc{Astuce}\ifx\relax#1\relax\relax\else\textmd{ : \itshape #1}\fi}}
\newtcolorbox{attention}[1][]{boxbase,colback=attncol!7,colframe=attncol,
  title={\textsc{Attention}\ifx\relax#1\relax\relax\else\textmd{ : \itshape #1}\fi}}

% Boîte "exercice" et "corrigé" (titre libre passé en argument)
\newtcolorbox{exobox}[1]{boxbase,colback=primary!4,colframe=primary,
  before skip=8pt,after skip=8pt,title={#1}}
\newtcolorbox{corbox}[1]{boxbase,colback=excol!5,colframe=excol!85!black,
  before skip=8pt,after skip=8pt,title={#1}}

% ---------------- En-tête ----------------
\newcommand{\docpart}[1]{\def\thedocpart{#1}}
\docpart{}
\pagestyle{fancy}
\fancyhf{}
\fancyhead[L]{\small\textcolor{primarydark}{\textbf{Fiche 8} — Statistiques descriptives}}
\fancyhead[R]{\small\textcolor{primarydark}{\textsc{\thedocpart}}}
\fancyfoot[C]{\small\thepage}
\renewcommand{\headrulewidth}{0.8pt}
\renewcommand{\headrule}{\hbox to\headwidth{\color{primary}\leaders\hrule height \headrulewidth\hfill}}

\titleformat{\section}{\large\bfseries\color{primarydark}}{\thesection}{0.6em}{}
\titleformat{\subsection}{\normalsize\bfseries\color{primary}}{\thesubsection}{0.6em}{}
\setlist{topsep=2pt,itemsep=1.5pt,parsep=0pt}

% Bandeau de titre du document
\newcommand{\bandeau}[2]{%
\begin{tcolorbox}[boxbase,colback=primary,colframe=primarydark,halign=center,
   before skip=2pt,after skip=12pt]
{\LARGE\bfseries\color{white} #1}\\[3pt]{\normalsize\color{white} #2}
\end{tcolorbox}}

\newcommand{\ecm}{\sigma_m}
\newcommand{\ect}{\sigma}
\newcommand{\med}{M\!e}
\docpart{Moyen — Thème 5}
\begin{document}
\bandeau{Thème 5 — Représentations graphiques}{Niveau Moyen — exercices}

\noindent\textit{Objectif : construire les angles d'un secteur, lire quartiles et boîte à moustaches.}\medskip

\begin{exobox}{Exercice 1 — Angles d'un diagramme circulaire}
Répartition des groupes sanguins de $N=300$ étudiants :
O : $45\,\%$ ; A : $40\,\%$ ; B : $10\,\%$ ; AB : $5\,\%$.
Calculer l'angle $\theta_i=f_i\times360^\circ$ de chaque secteur et vérifier que leur somme fait $360^\circ$.
\end{exobox}

\begin{exobox}{Exercice 2 — Lire $Q_1$, $\med$ et $Q_3$ sur l'ogive}
\begin{center}
\begin{tikzpicture}[scale=0.85]
  \draw[->,>=Stealth] (0,0)--(6.6,0) node[below right,font=\scriptsize]{valeur};
  \draw[->,>=Stealth] (0,0)--(0,4.5) node[left,font=\scriptsize]{$f_c$};
  \foreach \y/\v in {1/0{,}25,2/0{,}50,3/0{,}75,4/1{,}00}{\draw[black!12] (0,\y)--(6.2,\y);
     \draw (0,\y)--(-0.06,\y) node[left,font=\tiny]{\v};}
  \foreach \x/\lab in {1/10,2/20,3/30,4/40,5/50}\node[below,font=\tiny] at (\x,0){\lab};
  \draw[ogivecol,line width=1.1pt] (0,0)--(1,0.4)--(2,1.2)--(3,2.4)--(4,3.4)--(5,4);
  \foreach \x/\y in {1/0.4,2/1.2,3/2.4,4/3.4,5/4}\fill[ogivecol] (\x,\y) circle (1.4pt);
  \foreach \y in {1,2,3}\draw[medcol,dashed] (0,\y)--(6.2,\y);
\end{tikzpicture}
\end{center}
Les trois pointillés sont à $0{,}25$, $0{,}50$ et $0{,}75$. Entre quelles valeurs se situent
$Q_1$, $\med$ et $Q_3$ ?
\end{exobox}

\begin{exobox}{Exercice 3 — Lire une boîte à moustaches}
\begin{center}
\begin{tikzpicture}[scale=1]
  \draw[->,>=Stealth] (0,0)--(9.4,0) node[right,font=\scriptsize]{valeur};
  \foreach \x/\v in {1.6/4,2.8/7,4.4/11,6.4/16,8/20}
     \draw (\x,0.08)--(\x,-0.08) node[below,font=\scriptsize]{\v};
  \draw[line width=0.9pt] (1.6,0.7)--(2.8,0.7);
  \draw[line width=0.9pt] (6.4,0.7)--(8,0.7);
  \draw[line width=0.9pt] (1.6,0.45)--(1.6,0.95);
  \draw[line width=0.9pt] (8,0.45)--(8,0.95);
  \fill[primary!20,draw=primary,line width=0.9pt] (2.8,0.3) rectangle (6.4,1.1);
  \draw[medcol,line width=1.4pt] (4.4,0.3)--(4.4,1.1);
\end{tikzpicture}
\end{center}
Lire les cinq valeurs $x_{\min},Q_1,\med,Q_3,x_{\max}$, puis calculer l'écart interquartile.
\end{exobox}

\begin{exobox}{Exercice 4 — Type, graphique et mode}
Nombre de frères et sœurs des élèves d'une classe :
\begin{center}
\begin{tabular}{|c|c|c|c|c|c|}
\hline
$x_i$ & 0 & 1 & 2 & 3 & 4 \\\hline
$n_i$ & 5 & 11 & 8 & 4 & 2 \\\hline
\end{tabular}
\end{center}
(a) Le caractère est-il discret ou continu, et quel graphique lui convient ?
(b) Quel est le mode ? (c) Quelle fréquence (en \%) pour la valeur $1$ ?
\end{exobox}

\begin{exobox}{Exercice 5 — Proportion sur un diagramme circulaire}
Sur le diagramme circulaire des moyens de transport de $N=400$ étudiants, le secteur « Voiture »
occupe un angle de $90^\circ$. Combien d'étudiants viennent en voiture ?
\end{exobox}
\end{document}
